\section{Encoder、Decoder 与三种模型家族}

原始 Transformer 为机器翻译设计，包含 encoder 与 decoder。两者都堆叠 attention 和 FFN block，但信息访问不同。\term{encoder}（编码器）通常允许每个源 token 双向读取源序列；\term{decoder}（解码器）生成第 $t$ 个目标 token 时只能读取目标前缀，并通过 cross-attention 读取 encoder outputs。

\subsection{Encoder-decoder 数据流}

训练时，source tokens 经过 encoder 得到 memory；target tokens 右移后进入 masked decoder。Decoder self-attention 处理已生成前缀，\term{cross-attention} 用 decoder state 作为 query、encoder memory 作为 key/value，最后 softmax 预测下一个 token。\term{autoregressive}（自回归）表示联合概率分解为
\[
p(y\mid x)=\prod_{t=1}^{m}p(y_t\mid y_{<t},x).
\]

\lecturefigure{slide-14-encoder-decoder.jpg}{原始 Transformer 用 encoder memory、masked decoder 与 cross-attention 完成 seq2seq。}{官方视频 13:40；Vaswani et al. 2017。}

\begin{knowledgebox}{读图：三种 attention 不要混在一起}
Encoder self-attention 读取完整 source；decoder masked self-attention 只读取 target prefix；cross-attention 让 target query 读取 encoder memory。图中右移 target 确保模型不能直接看到当前答案。训练可并行计算所有 target position，但每个 position 的 mask 保持自回归语义；推理仍需逐 token 或使用专门并行解码方法。
\end{knowledgebox}

\begin{importantbox}{三种主流家族}
\begin{enumerate}
  \item \textbf{Encoder-only}：双向表示，适合分类、检索、标注与 representation learning；代表例为 BERT。
  \item \textbf{Decoder-only}：causal next-token modeling，适合开放式生成与统一 prompt interface；代表例为 GPT 系列。
  \item \textbf{Encoder-decoder}：source 与 target 分离，适合翻译、摘要和条件生成；代表例为原始 Transformer、T5。
\end{enumerate}
家族选择首先是信息边界和任务接口选择，不只是参数量选择。
\end{importantbox}

\subsection{本章小结}

Encoder、decoder 和 cross-attention 定义了信息从哪里来、何时可见。GPT 与 BERT 的差异可以从这张原始架构图直接推导，而不应被记成两个孤立模型名。

